\section{System prompt of the Autoresearch agent}
\label{app:prompt}

The following is the full text of the agent's program file. It is reproduced
verbatim because it is the evidence for the three prompt-side configuration
inconsistencies discussed in Section~\ref{sec:res-v2}; the remaining item, the
un-injected \texttt{total\_trials}, is a defect in the calling code and therefore
does not appear in this text. The annotated points are where the problems arise.

\begin{quote}
\footnotesize
\begin{verbatim}
# Autonomous HPO Agent - System Prompt

You are an autonomous hyperparameter optimization agent for medical VQA
fine-tuning research.

## Task
Given the history of previous QLoRA fine-tuning experiments, suggest the NEXT
hyperparameter configuration that is most likely to improve validation accuracy
on the PathVQA medical VQA dataset.

## Search Space
| Parameter        | Range                       | Type                     |
|------------------|-----------------------------|--------------------------|
| lora_rank        | {4, 8, 16, 32, 64}          | discrete                 |
| lora_alpha       | rank x {1, 2, 4}            | discrete                 |
| learning_rate    | [1e-5, 5e-4]                | continuous (log-scale)   |
| batch_size       | {1, 2, 4}                   | discrete                 |
| grad_accum_steps | {4, 8, 16}                  | discrete                 |
| warmup_ratio     | [0.0, 0.1]                  | continuous               |
| weight_decay     | [0.0, 0.1]                  | continuous               |
| lora_targets     | {"minimal","medium","full"} | categorical              |
| epochs           | {1, 2, 3, 5}                | discrete   <- (i) inert  |

Where: minimal = [q_proj, v_proj] / medium = [q_proj, k_proj, v_proj, o_proj]
       full = all linear layers

## Strategy Guidelines
1. Early exploration (trials 0-5): Try diverse configurations to map the
   landscape. Vary multiple parameters at once.
2. Mid exploitation (trials 5-20): Focus on promising regions. Take the best
   configuration and vary 1-2 parameters at a time.     <- (ii) budget mismatch
3. Late refinement (trials 20+): Fine-tune around the best configuration with
   small perturbations.                                 <- (ii) never triggered

## Key Insights for Medical VQA
- Medical images benefit from higher LoRA ranks (16-64).
- Learning rate is often the most sensitive parameter.
- `medium` or `full` target modules often outperform `minimal`.
- Effective batch size = batch_size x grad_accum_steps. Keep in 4-16 range.
- More epochs (3-5) help when training data is limited.  <- (i) inert parameter
- Warmup ratio 0.03-0.06 is generally safe.

## Response Format
Respond with ONLY a valid JSON object. No explanation, no markdown fences,
no other text.                                    <- (iii) rationale prohibited
\end{verbatim}
\end{quote}

\noindent
\textbf{(i) Inert parameter}: \texttt{epochs} appears in both the search space and
the guidance, yet the implementation discards the proposed value and fixes
\texttt{max\_steps}~$=200$.
\textbf{(ii) Budget mismatch}: the actual budget is 20 trials per repeat, so
``Late refinement (trials 20+)'' never triggers, and about 75\% of the budget
falls in the interval instructing the agent to take the best configuration and
vary 1--2 parameters.
\textbf{(iii) Rationale prohibited}: the prompt explicitly excludes the
interpretable search rationale required by RQ3.

\paragraph{Agent execution settings.}
Model \texttt{claude-sonnet-4-6}, \texttt{max\_tokens=512}. The temperature was
not a fixed value but a schedule that decreases with progress; an exhaustive tally
of the execution logs (432 calls) shows it ranged from 1.0 down to 0.66. The 0
recorded in the \texttt{temperature} column of the results table is an unpopulated
default and not the value actually passed to the API. Mitigation of API
non-determinism rests on the 10 independent repeats.

\section{Excerpts from the Autoresearch rationale logs}
\label{app:rationale}

Excerpted from the original text of the \texttt{agent\_reasoning} column of the
results table. Of the 200 completed trials, 147 (73.5\%) returned only JSON as in
(1) below, while 53 (26.5\%) included natural-language narrative.

\paragraph{(1) Typical response --- configuration only, no rationale (73.5\% of
the total).}

\begin{quote}
\footnotesize
\begin{verbatim}
{"lora_rank": 64, "lora_alpha": 256, "learning_rate": 2.0e-4, "batch_size": 2,
 "grad_accum_steps": 8, "warmup_ratio": 0.05, "weight_decay": 0.01,
 "lora_targets": "full", "epochs": 3}
\end{verbatim}
\end{quote}

\noindent
Because the prompt prohibited explanation (Appendix~\ref{app:prompt}, item (iii)),
this form is the response consistent with the instruction.

\paragraph{(2) A case including natural-language narrative --- accurate diagnosis,
but an inoperative lever.}
The agent diagnosed that all trials have only 200 steps, so this looks like
undertraining, and that epochs were not changed. The diagnosis was accurate, but
as noted in item (i) above the \texttt{epochs} lever it proposed was discarded by
the implementation and therefore never took effect.

\paragraph{(3) Repeated proposal of an identical configuration.}
In repeat 8, from trial 602 onward the agent proposed the identical configuration
11 times in succession, while the measured \texttt{val\_accuracy} varied between
0.388 and 0.444. This variation reflects the stochastic behaviour of training and
evaluation rather than any difference in configuration
(Sections~\ref{sec:res-behaviour}, \ref{sec:discussion}).

\section{Supplementary tables}
\label{app:tables}

\begin{table}[h]
\centering
\caption{PathVQA accuracy by clinical question type (zero-shot, seed 42).}
\label{tab:apptypes}
\small
\begin{tabular}{@{}lrrrrr@{}}
\toprule
Type & Samples & Gemma4-E2B & Qwen2.5-VL-3B & Qwen3-VL-2B & SmolVLM2-2.2B \\
\midrule
diagnosis   & 23    & 0.0000 & 0.0000 & 0.0000 & 0.0000 \\
location    & 433   & 0.0647 & 0.1062 & 0.1478 & 0.0762 \\
measurement & 33    & 0.2121 & 0.0909 & 0.1515 & 0.0909 \\
description & 2,729 & 0.0425 & 0.0224 & 0.0447 & 0.0198 \\
temporal    & 9     & 0.0000 & 0.0000 & 0.0000 & 0.0000 \\
yes\_no     & 3,362 & 0.1633 & 0.6130 & 0.6336 & 0.5892 \\
unknown     & 130   & 0.0692 & 0.0692 & 0.0923 & 0.0154 \\
\bottomrule
\end{tabular}
\end{table}

\begin{table}[h]
\centering
\caption{Decomposition of the accuracy gain by clinical question type
(Qwen3-VL-2B, PathVQA, 500 evaluation items). The best configuration is the best
trial of Autoresearch-v2. Share of gain is (type $n \div 500$) $\times$ per-type
gain, divided by the sum of positive contributions; the per-type contributions sum
to the overall gain of $+0.1260$. The same pattern appears under the best
configuration of the four-strategy comparison (yes\_no accounting for
53.4--63.2\% of the gain) and persists in the mean over the top ten trials of each
condition.}
\label{tab:appgain}
\small
\begin{tabular}{@{}lrrrrrr@{}}
\toprule
Type & WCA weight & $n$ & Zero-shot & Best config & Gain & Share of gain \\
\midrule
diagnosis   & 1.0 & 3   & 0.0000 & 0.0000 & $+0.0000$ & 0.0\% \\
location    & 0.8 & 21  & 0.2381 & 0.7143 & $+0.4762$ & 15.4\% \\
measurement & 0.7 & 4   & 0.0000 & 0.0000 & $+0.0000$ & 0.0\% \\
description & 0.6 & 211 & 0.0190 & 0.1137 & $+0.0948$ & 30.8\% \\
temporal    & 0.5 & 0   & ---    & ---    & ---       & --- \\
yes\_no     & 0.5 & 244 & 0.6721 & 0.8156 & $+0.1434$ & 53.8\% \\
unknown     & 0.5 & 17  & 0.1765 & 0.0588 & $-0.1176$ & $-3.1$\% \\
\midrule
Overall     & --- & 500 & 0.3520 & 0.4780 & $+0.1260$ & 100\% \\
\bottomrule
\end{tabular}
\end{table}

\begin{table}[h]
\centering
\caption{Indirect comparison with figures reported in prior work (closed-ended
basis). The LLaVA and LLaVA-Med figures are quoted from the original paper; the
figures for this study are \texttt{closed\_acc} for the Phase~2 main conditions
(mean of 3 seeds). The evaluation samples are identical at 1,061 items for SLAKE
and 451 for VQA-RAD; only PathVQA differs, 6,761 versus 6,719, a gap of 42 items.
Open-ended figures are excluded because the measures differ --- for reference,
LLaVA-Med's open figures are PathVQA 37.95 / SLAKE 83.08 / VQA-RAD 61.52 on a
token-recall basis, while ours are 17.15 / 66.95 / 26.00 on a
BERTScore-threshold basis.}
\label{tab:appllavamed}
\small
\begin{tabular}{@{}llrrr@{}}
\toprule
Model & Scale \& training approach & PathVQA & SLAKE & VQA-RAD \\
\midrule
LLaVA (general) \citep{liu2023llava} & 7B, no domain training & 63.20 & 63.22 & 65.07 \\
LLaVA-Med \citep{li2023llavamed} & 7B, biomedical pretraining $+$ per-dataset FT & 91.21 & 85.34 & 84.19 \\
This paper (Qwen3-VL-2B) & 2B, QLoRA fine-tuning only & 83.12 & 85.26 & 72.91 \\
\bottomrule
\end{tabular}
\end{table}
